\chap Performance

The performance of software models of \nns\ is probably their biggest
drawback ---
it is a classic tradeoff between flexibility and speed.
A hardware model is extremely inflexible, limited in size, and very 
cumbersome, yet it
offers performance that doesn't degrade as the network gets bigger.
Software models, on the other hand, are extremely flexible but
horrendously slow in {\sl all} implementations.
Here we examine just how slow a software simulation is.

An analysis of both execution and learning time, and the amount of
memory required for any network is given in the following brief
chapter.
The aim is to give the reader some simple rules of thumb for
calculating: how long a simulation using fishNET will take in any
particular implementation, and how big that simulation can be.

Problems with performance of fishNET tend to take the form of
execution speed, not lack of memory.
Long before available memory with a PC runs short, patience
inevitably does.
Lack of memory has actually only happened once, and this was a
deliberate experiment to see how large a network could be.
In reality, the time taken to learn (in physical seconds, not
internal learning time~$t$) becomes so large that it is nearly
pointless continuing.
Just how large is feasible?
Hopefully with the help of this chapter the reader can calculate how
long it will take and maybe rethink network size and shape.
Thus it is execution speed, not memory size, which is the most
important constraint in PC implementations.

\sec Hardware

A description of all the hardware configurations that fishNET was
tested with follows:

\ssec The (almost) Standard IBM-PC

As mentioned above, fishNET's development was undertaken entirely on
a 1985 vintage true-blue IBM-PC.
The~$4.77 {\rm MHz}$ int\lower.5ex\hbox{e}l~$8088$ has been replaced by a
National~${\rm V}20$ workalike processor, which is marginally faster
with screen and disk I/O and up to~$5$ times faster for some
processor operations.
The speed improvement, however, is only about~$20\%$ in general.
The machine has~$640{\rm k}$ of RAM.

\ssec The Turbo-Charging 8087

After finding that learning times for small networks was several
hours, the IBM-PC above was
fitted with a genuine int\lower.5ex\hbox{e}l~$8087$.
The speed improvement was in the order of a factor of~$12$.

\ssec The Pyramid Super-Mini

This machine literally turned days into minutes.
It is a Pyramid model~9810, which has~$16$~Mbytes of 
actual and~$200$~Mbytes of virtual
memory with a RISC architecture, so fishNET simulations 
of painful size for the PC were not much of a burden. 
Really.

\sec Learning and Execution Speed

To calculate learning and execution speed is easy on a machine as
slow as a PC, since it is possible to simply use a stop-watch and
observe the appropriate screen output.
On larger multi-tasking systems with buffered screen output this is
not possible.
Hence the timings shown for the Pyramid err on the high side
(the machine actually runs faster).
This is because it was necessary to use operating system information
on execution time which involved the complete running of the
program from invocation to termination, not just the different
phases such as learing and execution.
However, the case chosen for calculations had a reasonably large
number of learning iterations, so the result is accurate to within a
factor of~$1.5$.

\ssec Learning Speed

For simplicity's sake, consider the application of the
back-propagation algorithm to a weight as a single operation.
Then, it is possible to calculate the learning speed by measuring
the length of time taken for a single application of the
back-propagation algorithm and the application of $\Delta w(t)$,
then dividing by the number of I/O cases and dividing by the number
of weights.
(For more accuracy, several dozen sweeps were timed and the result
divided by the number of sweeps.)

This technique gave the following results:
$$
\vbox{\offinterlineskip
\hrule
\halign{&\vrule#&\strut\quad#\hfil\quad\cr
height2pt&\omit&&\omit&\cr
&Implementation&&
	Length of time for a `single' b.p. weight calculation&\cr
height2pt&\omit&&\omit&\cr
\noalign{\hrule}
height2pt&\omit&&\omit&\cr
&PC$_{\rm V20\ (8088)}$&&$9.17$\ msec&\cr
&PC$_{\rm 8087}$&&$733$\ $\mu$sec\ $(=0.733{\rm \ msec})$&\cr
&Pyramid&&$20$\ $\mu$sec\ $(=0.02{\rm \ msec})$&\cr
height2pt&\omit&&\omit&\cr}
\hrule}
$$
%$$\halign{#\hfil&\qquad#\hfil\cr
%Implementation&Length of time for `single' b.p. calculation\cr
%\noalign{\hrule}
%PC--V20		&$9.17 {\rm msec}$\cr
%PC--8087	&$733 \mu{\rm sec} (= 0.733 {\rm msec})$\cr
%Pyramid		&$20 \mu{\rm sec} (= 0.02 {\rm msec})$\cr
%}$$

A good approximation of speeds in general then is
$$
{\rm PC}_{\rm V20\ (8088)} = {{\rm PC}_{\rm 8087} \over 12.5} 
			= {{\rm Pyramid} \over 459}
$$

\ssec Execution Speed

Calculation of execution speed is a more difficult task.
Firstly, as after each execution there is a lot of disk output, the
stop-watch time for the PC will be grossly distorted (floppy disk
systems are even worse).
Hence the major mathematical operations involved will be analysed,
with their timings based upon reference [27].
Secondly, disk output time will be an order of magnitude slower than
processing speed in most implementations.
No attempt to analyse this time is made here, as it is too
application dependent.

For the simulator to operate upon a network (that is, use the
runtime data), there are a number of mathematical operations
involved. A summary is
\item{1.} $1\times$ transfer function calculation per \neu, which is
\itemitem{a)} $1 \times$ negate,
\itemitem{b)} $1\times$ {\tt exp} calculation,
\itemitem{c)} $1\times$ {\tt double} add, and
\itemitem{d)} $1\times$ {\tt double} divide.
\item{2.} $1\times$ {\tt double} multiplication per weight, and
\item{3.} $1\times$ {\tt double} addition per weight.

\noindent
Hence, using [27], we find that the approximate execution speed for
a single input case is of the order of
$$
\eqalign{
{\rm PC}_{\rm V20\ (8088)} &= (n_{\rm \neus} \times 41.61 \times 10^{-3})\cr
	&\qquad{}+ (n_{\rm weights} \times (0.753\times 10^{-3} + 
				2.71\times 10^{-3}))
				{\rm\ sec}\cr
		   &= (n_{\rm \neus} \times 41.61) + 
	(n_{\rm weights} \times 3.463) {\rm \ msec}\cr
{\rm PC}_{\rm 8087} &= (n_{\rm \neus} \times 1.08\times 10^{-3})\cr
	&\qquad{}+ (n_{\rm weights} \times (0.186\times 10^{-3} + 0.187\times
		10^{-3})) {\rm\ sec}\cr
	      &= (n_{\rm \neus} \times 1.08) +
	(n_{\rm weights} \times 0.373) {\rm \ msec}\cr
}$$
These values are very approximate.
However, they still show an interesting trend.
Assuming that the Pyramid is at least~$37$ times 
faster than the PC$_{\rm 8087}$, we find
$$
\eqalign{
{\rm Pyramid} &= (n_{\rm \neus} \times 29.2) +
		(n_{\rm weights} \times 10.1){\rm\ }\mu{\rm sec}\cr
}
$$
So for a typical problem with~$3$ layers, ($195$ \neus\ in the input
layer, $30$ in the middle layer, and $5$ in the output layer, or
$6000$ weights in total,) the expected execution times are
$$
\eqalign{
{\rm PC}_{\rm V20\ (8088)} &= 30.35 {\rm\ sec}\cr
{\rm PC}_{\rm 8087}	    &= 2.49  {\rm\ sec}\cr
{\rm Pyramid}	    &= 67    {\rm\ msec}\cr
}
$$
These are close to the roughly measured times.

As you can see, the fact that there are so many more weights than
\neus\ means that calculations involving weights consume the most
significant part of the time, even though the execution speed of a
\neu\ is typically~$3 \to 12$ times slower.
In these calculations, memory addressing time has been ignored, as
it is~$100 \to 1000$ times faster than floating point operations.

\sec Memory Used

Once again, we are only searching for a rule of thumb: on single
processor systems such as PCs, memory resident device drivers and
RAM disks can make it difficult to accurately guage free memory
(the MS--DOS utility {\tt chkdsk} gives this information, as well as
disk free data), and on
multi-tasking systems, reasonably sized simulations aren't going to
pose much of a problem.

A quick calculation follows.
Memory used is broken into four categories:
\item{1.} a structure of type {\tt QUESTION},
\item{2.} the network itself,
\item{3.} a variable amount of space for teaching and execute data, and
\item{4.} a variable amount of space for calculation purposes.

The space taken up by the structure of type {\tt QUESTION} is
intentionally always the same.
The actual number of bytes varies slightly from implementation to
implementation (as the size of an {\tt int} varies).
For a PC application
$$
\eqalign{
{\tt QUESTION}_{\rm space} &= 4 \times ( 36 \times {\tt char} +
	2 \times {\tt int}) + (10 \times {\tt int}) \cr
		&\qquad + (2 \times {\tt double}) +
			(114 \times {\tt char})\cr
		&= (258 \times {\tt char}) + (2 \times {\tt double}) +
			(18 \times {\tt int})\cr
		&= (258 \times 1) + (2 \times 8) + (18 \times 2)\cr
		&= 310 {\rm\  bytes}.
}
$$
It will become obvious that this is hardly significant.

The network requires space for
\item{1.} {\tt LAYER} structures,
\item{2.} {\tt NEURON} structures, and
\item{3.} {\tt WEIGHT} structures.

Hence the amount of space required is
$$
\eqalign{
{\tt LAYER}_{\rm space} &= ({\tt NEURON *}) = 4 {\rm\ bytes},\cr
{\tt NEURON}_{\rm space}&= ({\tt WEIGHT *}) + {\tt double} 
			= 12 {\rm \ bytes}, \cr
{\tt WEIGHT}_{\rm space}&= 3 \times {\tt double} = 24 {\rm \ bytes}.\cr
}
$$
So the equation for memory usage by a network becomes
$$
{\tt NETWORK}_{\rm space} = (n_{\rm layer} \times 4) + 
	(n_{\rm neuron} \times 12) + (n_{\rm weight} \times 24),
$$
where~$n_{\rm layer}$ is the number or layers, etc.

Memory required for teaching and execution data depends upon
the number of \neus\ in the input and output layers, and the
number of I/O cases.
Every \neu\ on the input has a piece of data for learning and execution
for each case.
Every \neu\ on the output has a piece of data for comparison during learning
for each I/O case.
So
$$
\eqalign{
{\rm data}_{\rm space} &= (n_{\rm layer\ 0\ \neus} \times 
	({\rm learning\ cases} + {\rm execution\ cases}) \times 
	{\tt double})\cr
		&\qquad{}+ (n_{\rm output\ layer\ \neus} \times
		{\rm learning\ cases}\times {\tt double})\cr
		       &= (n_{\rm layer\ 0\ \neus} \times 
		({\rm learning\ cases} + {\rm execution\ cases}) \times 8)\cr
			&\qquad{}+ (n_{\rm output\ layer\ \neus}\times
				{\rm learning\ cases} \times 8).\cr
}
$$

The next facet to consider is space used during calculations.
Arrays of type {\tt double} are used during back-propagation
calculations (in the functions {\tt back\-\_propagate} and {\tt
back\-\_propagate\-\_apply\-\_delta\-\_w} as described in the
previous chapter) to store values of $\partial E \over \partial x$
and $\partial E \over \partial y$.
The arrays refer to different layers of \neus\ at a time, but for
the rule of thumb we will invent a worst case.
The worst case is where we have two big layers of the same size
next to each other, so
$$
{\rm calculation}_{\rm space} = (n_{\rm \neus\ in\ biggest\ layer}\times 16)
$$

Putting all these parts together, 
$$
\eqalign{
{\rm Memory\ required} &= {\tt QUESTION}_{\rm space} + 
		{\tt NETWORK}_{\rm space} + {\rm data}_{\rm space}
		+ {\rm calculation}_{\rm space}\cr
&= 310 + (n_{\rm layer} \times 4) +
	(n_{\rm \neu} \times 12) + (n_{\rm weight} \times 24)\cr
	&\qquad{}+(n_{\rm layer\ 0\ \neus} \times ({\rm learning\ cases}
		+{\rm execution\ cases}) \times 8)\cr
	&\qquad{}+(n_{\rm output\ layer\ \neus} \times 
		{\rm learning\ cases} \times 8)\cr
	&\qquad{}+(n_{\rm \neus\ in\ biggest\ layer} \times 16) {\rm\ bytes}.\cr
}
$$
Using again the example of a 3 layer (195, 30, 5) network, (with 5 learning
cases and 20 execution cases,) we
get
$$
\eqalign{
{\rm Memory\ required} &= 310 + (3 \times 4) + (230 \times 12) + 
			(6000 \times 24)\cr
			&\qquad{}+ (195 \times 25 \times 8) +
			(5 \times 20 \times 8) + (195 \times 16)\cr
		      &= 190\,002 {\rm\ bytes}\cr
		      &= 185.5 {\rm\ kbytes}.
}
$$

\sec An Interesting Comparison --- Software versus Wetware

The human brain has been dubbed `wetware' by some authors.
How does wetware compare with a single processor software simulation
of the brain?
Consider a very, very simple example.

According to Feldman et.\ al.\ [12], the ``execution'' time for a
human \neu\ to produce an output from its inputs is in the order
of~$5$ milliseconds.
However, {\sl all} \neus\ work at the same time (within a layer,
say) so a~$3$ layer network should produce an output in~$15$
milliseconds.
Considering the gross simplifications here, let us increase this
estimate by an order of magnitude, and say that for a simple
perceptual problem, wetware produces an answer in~$200$
milliseconds.
Assume we used a tiny fraction of the brain, say~$300\,000$ \neus.
Assume again that we have only~$5\,000$ connections per \neu.
Then if the~$300\,000$ are equally divided in 3 layers, we
have~$1\times 10^9$ weights.
And it can produce an answer in~$200$ milliseconds.

A Pyramid would require $\sim 7.7 \times 10^3$ seconds $=$ $2$ hours
$9$ minutes to do the same calculation, and a plain PC would take
$\sim 41$ days $4$ hours.
This also involves {\sl no} consideration for how long it would take
to teach the software models by back-propagation.

Software artificial \nns\ will never be able to perform a 
meaningful portion
of the human brain's functions, no matter how fast the machine, if
the only available technology is a few planar silicon processors.
True multidimensional parallel processing such as optical or
biological processors offer the only possible path.
Or we will be stuck with simulations of the same order as the ones
investigated in this thesis.
Simulations of this size are however remarkably useful, as the next
part of this thesis shows.

